\subsection{An iterated PD saturation with actual target maps}

\begin{theorem}[A-SAT01: Proved; simultaneous saturation in all sufficiently large weights]
\label{ASAT01}
Let $R=\mathbb Z$ or $\mathbb Z_{(p)}$ and $S_R=\Gamma_R(T)$.
Let $X$ be a nonnegative weight-graded strict simplicial PD cochain
$S_R$-algebra with torsionfree terms. Suppose $X_0=S_R$,
$H^0(X_n)=S_R$, and $S_{\mathbb Q}\to X_{n,\mathbb Q}$ is a
quasi-isomorphism in every simplex and weight. Let
$f_i:X\to A_i$ be a specified set of strict maps to torsionfree
targets. Choose a positive integer $W$ such that, for every $i$
and $w\geq W$, the coefficient comparison gives
$H^0(A_{i,n}^{w,*})=S_R^w$ and $H^1(A_{i,n}^{w,*})=0$,
uniformly in $n$. At least one target is specified.

There is an actual strict inclusion $X\hookrightarrow X^{\mathrm{sat}}$
and extending strict maps $f_i^{\mathrm{sat}}:X^{\mathrm{sat}}\to A_i$
with the following properties:
\begin{enumerate}
\item all weights below $W$ are unchanged;
\item the terms are torsionfree, hence flat over $R$;
\item rationalization gives the same algebra $X_{\mathbb Q}$,
      and $H^0(X^{\mathrm{sat}}_n)=S_R$ in every weight;
\item $H^1((X^{\mathrm{sat}}_n)^{w,*})=0$ for $w\geq W$;
\item $X^{\mathrm{sat}}_0=X_0=S_R$;
\item when $S_R\to X_1$ is an integral quasi-isomorphism
      in every weight, $X^{\mathrm{sat}}_1=X_1$ literally.
\end{enumerate}
The support condition $q\leq w$ is preserved whenever imposed on $X$.
No preservation of horizontal exactness or of higher integral
vertical cohomology is part of this assertion.
\end{theorem}

\begin{proof}
We construct every stage inside the fixed rational PD algebra
$X_{\mathbb Q}$. Since the targets are torsionfree, we also identify
$A_i$ with its actual image in $A_{i,\mathbb Q}$.
For any intermediate algebra $D$ containing $X$ in $X_{\mathbb Q}$,
with the stated coefficient $H^0$, simplex-zero algebra, and maps
to the targets, define in each simplex and weight $w\geq W$
\[
 E_n^w(D)=\{a\in X_{n,\mathbb Q}^{w,0}:
              da\in D_n^{w,1},\quad a|_0\in S_R^w\}.
 \tag{SAT.1}
\]
Here $a|_j$ is restriction along the $j$th vertex; it lies in
$S_{\mathbb Q}^w$ because $D_0=X_0=S_R$.

All vertex values of an element of (SAT.1) are integral.
Indeed $d a$ is a closed integral one-form in $D_n$.
For a specified target, choose $h\in A_{i,n}^{w,0}$ with
$dh=f_i(da)$, using its vanishing $H^1$.
Each $h|_j$ is closed, since $(da)|_j=0$ in $D_0^1$.
It is therefore the integral coefficient in $S_R^w$ given by
the target's $H^0$ comparison. The rational difference
$f_{i,\mathbb Q}(a)-h$ is a closed zero-form, so it is a
single coefficient in $S_{\mathbb Q}^w$. Its values at all
vertices agree. Consequently
\[
                  a|_j-a|_0=h|_j-h|_0\in S_R^w.
 \tag{SAT.2}
\]
This proves the asserted integrality. It also proves that
$E(D)$ is stable under every simplicial operator, since a
new vertex-zero restriction is one of these integral vertex
values. It is an additive graded $S_R$-module: coefficients
are closed, and multiplication by a coefficient preserves
the condition on $da$ and on vertex values.

For each target the actual value of the extending map on $a$ is
\[
             \widetilde f_i(a)=h+(a|_0-h|_0)\in A_{i,n}^{w,0}.
 \tag{SAT.3}
\]
Two choices of $h$ differ by an integral coefficient, so
(SAT.3) is independent of that choice. It is also exactly
$f_{i,\mathbb Q}(a)$: their difference is a rational coefficient
whose vertex-zero value is zero. In particular it is additive,
$S_R$-linear, and simplicial. The same argument supplies these
values for every specified target, with no relation between
the different choices of primitives required.

Let $T(D)$ be the PD subalgebra of $X_{\mathbb Q}$ generated by
$D$ and every $E_n^w(D)$ for $w\geq W$, in all simplices.
It is simplicial by the preceding stability. It is a differential
subalgebra: the derivatives of the added zero-forms already
belong to $D$, and the derivation and PD differentiation rules
then preserve the generated subalgebra. All new generators have
positive weight, so the augmentations are preserved and no weight
below $W$ changes. Formula (SAT.3) says that the restriction of
each actual rationalized map $f_{i,\mathbb Q}$ takes these generators
to $A_i$. Products and divided powers therefore remain in $A_i$.
This gives strict extending maps on the entire $T(D)$, including
every relation, because they are restrictions of the already
defined rational algebra maps to torsionfree subalgebras.

The new terms are torsionfree. The inclusions
$X\subset D\subset T(D)\subset X_{\mathbb Q}$ give
$T(D)_{\mathbb Q}=X_{\mathbb Q}$.
Every closed zero-form in $T(D)$ is therefore an actual rational
coefficient, since the rational coefficient comparison has no
degree-minus-one boundaries. In a changed weight its image in
any specified $A_i$ is an integral coefficient by $H^0(A_i)=S_R$.
The rational map fixes coefficients, so the original coefficient
already lies in $S_R$. In a smaller weight the algebra was unchanged.
Thus $H^0(T(D)_n)=S_R$ in every weight.

At simplex zero, (SAT.1) adds only coefficients already in $S_R$;
hence $T(D)_0=S_R$. Suppose also $D_1=X_1$ and that $X_1$
has the integral coefficient comparison in all weights.
For $a\in E_1^w(D)$ the integral closed form $da$ has a primitive
$b\in X_1^{w,0}$. Its vertex-zero value is integral.
Replace $b$ by $b+(a|_0-b|_0)$.
The rational difference $a-b$ is then a closed zero-form
with vertex-zero value zero, hence is zero.
Therefore $a\in X_1$, and $T(D)_1=X_1$.

Starting with $D_0=X$, iterate $D_{r+1}=T(D_r)$ for
$r\geq0$ and put $X^{\mathrm{sat}}=\bigcup_{r\geq0}D_r$.
All inclusions and target maps are the actual restrictions
described above, so this union is a strict simplicial PD algebra
with the asserted compatible maps. The preceding arguments
preserve every stated invariant at every stage and in the union.

Finally, a closed one-form $c\in(X^{\mathrm{sat}}_n)^{w,1}$
with $w\geq W$ lies in some $D_r$. Its derivative is already
zero there because all stages are subalgebras of $X_{\mathbb Q}$.
Rational acyclicity gives $a\in X_{n,\mathbb Q}^{w,0}$ with $da=c$.
Subtract the rational coefficient $a|_0$ from $a$.
The new primitive has vertex-zero value zero and belongs
to $E_n^w(D_r)$, so it lies in $D_{r+1}$.
Thus $c$ is a boundary in the union, proving the required $H^1$
vanishing. The support assertion follows from containment in
the originally supported rational algebra. Flatness follows
from torsionfreeness over the PID $R$.
\end{proof}

In particular the theorem applies to a resolved-tail algebra
of AT16 when its original source has the rational coefficient
comparison. The targets $A,B$ are retained. An additional actual
comparison to a torsionfree two-term tail algebra with coefficient
$H^0$ and vanishing $H^1$ can be included in the specified target
set, using a common cutoff $W$. Formula (SAT.3) then extends that
very comparison through saturation. Thus this construction gives
an actual saturated comparison with $H^1=0$ in the stated range;
its effect on $H^2$ and $H^3$ remains a separate calculation.
